% !TEX root = userguide.tex

\def\headdate{28th March 2023}

\section{The generalized Stokes equation}
\label{sec:genstokes}

In this section some examples using the Stokes equation with
a varying coefficient (viscosity) $\eta$
will be given. We will refer to this as the generalized Stokes equation
in contrast to the usual Stokes equation, which has a constant coefficient $\eta$.
See Sec.~\ref{sec:stokesweak} for the theory, which is valid for both constant
and varying viscosity. For solving flows with generalized Newtonian fluids, we refer to
Sec.~\ref{sec:gennewtonian}.

\subsection{A generalized Stokes problem on a unit square with periodic boundary
         conditions: channel problem with specified flowrate; collocation;
         layers with different viscosity.}
\label{sec:genstokes1}
                                                                                
We consider the generalized Stokes equation for a channel problem on a unit square
similar to the problem \texttt{stokes2.f90} in Sec.~\ref{sec:stokes2}. Near the upper
and lower wall there exist fluid layers that have a lower viscosity than the layer in the middle. 
                                                                                
The example files are \texttt{generalized\_stokes1.f90} and \texttt{generalized\_stokes2.f90}
and can be obtained by typing at the
command line:
\begin{quote}
   \texttt{getexample generalized\_stokes}
\end{quote}
In both examples the viscosity is a function of $y$:
\[
     \eta(y) = 
\begin{cases}
 0.2 & \quad y\in[0,0.2)\text{ or }y\in(0.8,1]\\
  1  & \quad y\in(0.2,0.8)
\end{cases}
\]
In example \texttt{generalized\_stokes1.f90} the viscosity is computed
using a function routine in the
element. 
In \texttt{generalized\_stokes2.f90} the problem is solved
by filling a vector defined
per element of type \texttt{elvector\_t} with the Gauss point values of
$\eta$. The resulting velocity profile has been plotted in
Fig.~\ref{fig:vprofile1}.
\begin{figure}[htbp!]
   \begin{center}
   \includegraphics[scale=0.55]{vprofile1}
   \end{center}
   \caption{Velocity profile in the generalized\_stokes1 and 
    generalized\_stokes2 problem.}
    \label{fig:vprofile1}
\end{figure}

\subsection{A generalized Stokes problem on a unit square with periodic boundary
         conditions: channel problem with specified flowrate; collocation;
         Generalized Newtonian fluid model.}
\label{sec:genstokes3}
                                                                                
We consider the generalized Stokes equation for a channel problem on a unit square
similar to the problem \texttt{stokes2.f90} in Sec.~\ref{sec:stokes2}. The
viscosity is now given by a generalized Newtonian model.
                                                                                
The example files are \texttt{generalized\_stokes3.f90} and
\texttt{generalized\_stokes4.f90},
and can be obtained by typing at the
command line:
\begin{quote}
   \texttt{getexample generalized\_stokes}
\end{quote}
In \texttt{generalized\_stokes3.f90} the problem is solved
by filling a vector defined
per element of type \texttt{elvector\_t} with the Gauss point values of
$\eta$ according to a generalized Newtonian model. The solution is obtained
using a Picard iteration scheme. See also Sec.~\ref{sec:gennewtonian}, where this problem is solved using the more convenient module \texttt{generalized\_newtonian\_elements\_m}.

In \texttt{generalized\_stokes4.f90},
the viscosity is given by a regularized Bingham model \cite{Papanastasiou87}:
\[
     \eta(\dot\gamma) = \eta_0 +
       \dfrac{\tau_\text{y} \left( 1 - \exp(-r \dot\gamma) \right)}{\dot\gamma}
\]
where $\eta_0$ is the viscosity of the yielded material, $\tau_\text{y}$ is the yield stress, and $r$ is the regularization parameter.
The resulting velocity profile has been plotted in Fig. \ref{fig:vprofile3}.

\begin{figure}[htbp!]
   \begin{center}
   \includegraphics[scale=0.55]{vprofile3}
   \end{center}
   \caption{Velocity profile in the generalized\_stokes4 problem.}
    \label{fig:vprofile3}
\end{figure}
